
\subsubsection*{4.1.2}
\phantomsection
\label{sga3:VIIB:4.1.2}

Let \(G\) be a formal \(k\)-group. Since the functor
\[
  X\longmapsto X^{(p)}
  =
  X\widehat\otimes_\pi k
\]
commutes with finite products, it carries formal \(k\)-groups to
formal \(k\)-groups. Moreover, \(\operatorname{Fr}(X/k)\) is
functorial in \(X\), so
\[
  \operatorname{Fr}:G\longrightarrow G^{(p)}
\]
is a homomorphism of formal \(k\)-groups.

Define recursively
\[
  G^{(p^{\,n+1})}
  =
  \bigl(G^{(p^n)}\bigr)^{(p)}.
\]
Then the same is true of the composite
\[
\begin{aligned}
  \operatorname{Fr}^{n}
  =
  \operatorname{Fr}^{n}(G/k):
  G
  &\xrightarrow{\ \operatorname{Fr}\ }
  G^{(p)}
  \xrightarrow{\ \operatorname{Fr}\ }
  G^{(p^2)}
  \xrightarrow{\ \operatorname{Fr}\ }
  \cdots
  \xrightarrow{\ \operatorname{Fr}\ }
  G^{(p^n)}.
\end{aligned}
\]

Let \(A\) be the affine algebra of \(G\), and let \(I_A\) be its
augmentation ideal.

\medskip
\noindent\textbf{Definitions.}
\begin{enumerate}
\item[\textup{(a)}] The kernel of \(\operatorname{Fr}^{n}\) is
denoted
\[
  \operatorname{Fr}_nG.
\]
It follows directly from the definitions that
\(\operatorname{Fr}_nG\) is infinitesimal and has affine algebra
\[
  A/I_A^{\{p^n\}},
\]
where \(I_A^{\{p^n\}}\) is the closed ideal of \(A\) generated by
the \(p^n\)-th powers of the elements of \(I_A\).

\item[\textup{(b)}] The group \(G\) is said to have
\emph{height at most \(n\)} if
\[
  I_A^{\{p^n\}}=0;
\]
equivalently,
\[
  \operatorname{Fr}_nG=G.
\]
\end{enumerate}
